\section{Curve alternatives}
\label{sec:curves}

It remains to supply the one-dimensional Hilbert bounds for the affine
counting theorem.  A curve with a nonpolynomial rational coordinate will
instead contain very few selected points.  We first prove that assertion
uniformly, even when the rational function varies with the curve and
with the integral coordinate changes.

\begin{lemma}[Integral values of a rational function]\label{lem:rational-values}
Fix a positive integer $D$ and positive constants $A,B,C_0$.
For every $\eps>0$ and $X\ge2$, let $R\in\Q(s)$ be a nonpolynomial
rational function whose reduced numerator and denominator have degrees
at most $D$.  Then
\[
 \#\{s\in\Z: |s|\le C_0X^A,\ R(s)\in\Z,
                         \ |R(s)|\le C_0X^B\}
 \ll_{D,A,B,C_0,\eps}X^\eps.
\]
The constant is independent of the coefficients of $R$.  Inputs at
which $R$ has a pole are excluded from the set.
\end{lemma}

\begin{proof}
Write $R=A_0/B_0$ in lowest terms, with true degrees $a,b\le D$.
Since $R$ is not polynomial, $b\ge1$.  If there are at most $a+b$
counted inputs, the result is immediate.  Otherwise choose $a+b+1$
distinct input-value pairs $(s_i,v_i)$ and impose the homogeneous linear
equations
\[
 A_1(s_i)-v_i B_1(s_i)=0,\qquad
 \deg A_1\le a,\quad \deg B_1\le b.
\]
Their solution space is one-dimensional.  Indeed, for any solution
the polynomial $A_1B_0-A_0B_1$ has degree at most $a+b$ and vanishes
at all the chosen inputs.  It is therefore zero.  Coprimality of
$A_0,B_0$ and the true degree restrictions force
$(A_1,B_1)$ to be a scalar multiple of $(A_0,B_0)$.

The integer matrix of this system has polynomially bounded entries in
$X$.  Its maximal signed minors give a nonzero integer solution
$(A_1,B_1)$ with coefficients $X^{O_{D,A,B,C_0}(1)}$.  The resultant
\[
 R_0=\operatorname{Res}(A_1,B_1)
\]
is a nonzero integer of size $X^{O_{D,A,B,C_0}(1)}$.  At a counted
input, $B_1(s)$ is nonzero and divides $A_1(s)$.  The integral Bezout
identity for the resultant consequently gives $B_1(s)\mid R_0$.
The integer $R_0$ has $O_\eps(X^\eps)$ signed divisors, after adjusting
the exponent in the divisor bound.  For each such divisor $q$, the
nonconstant polynomial equation $B_1(s)=q$ has at most $b$ solutions.
This proves the uniform estimate.
\end{proof}

We apply this lemma to the selected arithmetic points from
Section~\ref{sec:arithmetic-boxes}, after the integral changes of
coordinates in Proposition~\ref{prop:lattice-boxes}.  Thus the
coordinates are $x',y'\in\Z^d$, with respective weights $U,V>0$ and
size bounds
\[
 |x'_i|\ll_f X^U,\qquad |y'_i|\ll_f X^V.
\]
The conjugate forms $a_i,b_i$ are complex linear combinations of $x'$
and $y'$ respectively.  Their filtered weights are at most $U$ and
$V$, regardless of the sizes of those linear coefficients.  They
satisfy $a_i^k b_i=n-\theta_i$.  The function
\begin{equation}\label{eq:curve-input-function}
 n=\frac1d\sum_{i=1}^d(a_i^k b_i+\theta_i)
\end{equation}
is defined over $\Q$: its coefficients are traces of elements of $K$.
It is a polynomial of degree at most $k+1$ in the rational coordinates.

Let $S$ denote any resulting selected set, with the global exceptional
values of Proposition~\ref{prop:quintic-discard} removed when $d=4$.
By Lemma~\ref{lem:balanced-factorization}, a fixed $n$ costs at most one
point of $S$, and a fixed $x'$ costs $O_f(1)$ points.  These bounds are
preserved when $S$ is restricted to any of the algebraic subsets used
in the counting argument.

\begin{proposition}[Curve alternatives]\label{prop:curve-alternatives}
Put
\[
 g=(d-1)(k-1)+\mathbf 1_{d=4},\qquad
 \delta_1(U,V)=\max\{U/2,\min(U,V/g)\}.
\]
Fix a degree bound $E$, and let $C$ be a geometrically irreducible
curve of degree at most $E$, defined over $\Q$, in the arithmetic
variety $a_i^k b_i=n-\theta_i$ in the coordinates just described.
For every $\eps>0$, either
\[
 |S\cap C|\ll_{f,E,\eps}X^\eps,
\]
or its filtered Hilbert dimension, with weight $U$ on $x'$ and weight
$V$ on $y'$, satisfies
\begin{equation}\label{eq:curve-hilbert-bound}
 H_C(T)\ge \frac{T}{\delta_1(U,V)}-O_{f,E}(1).
\end{equation}
The implied constants may also depend on a fixed compact subset of
$(0,\infty)$ containing $U,V$.  They are otherwise uniform in the
weights and in the integral coordinate changes.  In particular they
are uniform for the weights arising from
Proposition~\ref{prop:lattice-boxes}.
\end{proposition}

\begin{proof}
If $n$ is constant on $C$, the fixed-$n$ bound applies.  If the
$a$-image is constant, invertibility of the coordinate changes makes
$x'$ constant, so the fixed-$x'$ bound applies.  We henceforth assume
both are nonconstant.

Suppose first that the $a$-image is a nonlinear curve.  A generic
linear projection of that image to an affine line has degree equal
to its degree, hence at least two.  This follows by intersecting the
projective closure with general hyperplanes in the corresponding
pencil; the center can be chosen to avoid the points at infinity.
Let $t$ be the projected linear function, of weight at most $U$.
The field $\C(C)$ then has degree at least two over $\C(t)$.
When the $a$-image is a line, choose a nonconstant $x'$-coordinate as
$t$; the same conclusion holds if the function field of $C$ has
degree at least two over that line.
In either case some $x'$- or $y'$-coordinate $q$ lies outside $\C(t)$,
since the ambient coordinates generate $\C(C)$.  The families
\[
 1,t,t^2,\ldots,\qquad q,qt,qt^2,\ldots
\]
are linearly independent.  The weight of $q$ is at most $\max(U,V)$,
so counting these functions gives
\begin{equation}\label{eq:curve-two-families}
 H_C(T)\ge 2T/U-O(1).
\end{equation}
The error is uniform for the stated weights.  Since
$\delta_1(U,V)\ge U/2$, this proves
\eqref{eq:curve-hilbert-bound} in these cases.

It remains to treat a line image with function-field extension degree
one.  Work in the rational coordinates $x'$.  The image line is
defined over $\Q$, as the image closure of a $\Q$-defined curve under
a $\Q$-linear projection.  Choose a nonconstant coordinate
$s=x'_j$ as its parameter.  Each $x'_i$ is an affine rational function
of $s$.  Undoing the integral coordinate change and using the integral
basis of $K$, the line takes the form
\[
 \alpha=rs+q,\qquad r,q\in K,\quad r\ne0.
\]
Consequently
\begin{equation}\label{eq:curve-conjugate-line}
 a_i=r_i s+q_i,\qquad
 r_i=\sigma_i(r)\ne0,\quad q_i=\sigma_i(q),\quad r_i,q_i\in L.
\end{equation}
Every conjugate slope is nonzero because a field embedding cannot
annihilate a nonzero element.

The degree-one field extension gives $\C(C)=\C(s)$.  Each $y'_i$, and
the function $n$ in \eqref{eq:curve-input-function}, therefore belongs
to $\Q(s)$.  Indeed its expression as a rational function of the
$\Q$-defined parameter $s$ is unique when written with a coprime
numerator and monic denominator.  All coefficients are therefore fixed
by automorphisms over $\Q$.  These rational functions have uniformly
bounded degrees.
For $y'_i$, project the curve to the $(s,y'_i)$-plane: its graph has
degree at most $E$, and its reduced equation is
$B(s)y'_i-A(s)=0$.  The same argument with the fixed-degree map
\eqref{eq:curve-input-function} bounds the degree for $n$.

If any $y'_i$ is a nonpolynomial rational function, apply
Lemma~\ref{lem:rational-values} to the integral inputs $s$ and the
integral outputs $y'_i$.  Their bounds follow from the transformed
coordinate box.  Each $s$ fixes the point $x'$ on the line, so there
are only $O_f(1)$ selected points for that input.  This proves the
exceptional counting alternative.  This argument requires no height
bound on the coefficients of the moving line or coordinate change.

We may now assume that every $y'_i$ is polynomial in $s$.  Then the
$b_i$ and $n$ are polynomials as well.  Put $\tau_i=-q_i/r_i$.
The identities give
\[
 (s-\tau_i)^k\mid n(s)-\theta_i.
\]
The numbers $\tau_i$ are distinct: a common value would force two
distinct conjugates $\theta_i$ to equal $n$ at that value.  As $n$ is
nonconstant, differentiation yields
\begin{equation}\label{eq:curve-polynomial-degrees}
 \deg n\ge1+d(k-1),\qquad
 \deg b_i=\deg n-k\ge(d-1)(k-1).
\end{equation}
The first inequality follows because
$\prod_i(s-\tau_i)^{k-1}$ divides $n'(s)$; the degree identity uses
the nonzero slope of each $a_i$.

For $d=4$, equality in the first inequality would give $\deg n=5$.
Its derivative would then have precisely the four simple zeros
$\tau_i$, with respective critical values $\theta_i$.  Define
\[
 \Phi(z)=n\bigl(\tau_1+(\tau_2-\tau_1)z\bigr).
\]
This polynomial belongs to $L[z]$, has degree five, and has normalized
critical points $0,1$ and two further distinct points with the
prescribed ordered critical values.  At a selected point, the integer
$s$ gives $z=(s-\tau_1)/(\tau_2-\tau_1)\in L$.  Thus its integer value
$n=\Phi(z)$ belongs to the global exceptional set of
Proposition~\ref{prop:quintic-discard}.  Such a curve has no surviving
selected points.  On every remaining polynomial line curve,
\[
 \deg b_i\ge g=(d-1)(k-1)+\mathbf 1_{d=4}.
\]
The exceptional set was removed once, before choosing curves; its
bound is not charged separately to different lines or boxes.

Finally choose one $b_i$, of degree $\ell\ge g$ in $s$.  The functions
\[
 s^j b_i^a,\qquad 0\le j<g,\quad a\ge0,\quad jU+aV\le T,
\]
are linearly independent, since their degrees $j+\ell a$ as
polynomials in $s$ are distinct.  Their weights are at most $jU+aV$.
They therefore give
\[
 H_C(T)\ge gT/V-O_g(1+U/V).
\]
Powers of $s$ alone also give $H_C(T)\ge T/U-O(1)$.  The ratio $U/V$
is bounded in the stated range, so these two estimates imply
\[
 H_C(T)\ge
 \frac{T}{\min(U,V/g)}-O(1).
\]
Together with \eqref{eq:curve-two-families}, this proves the proposition.
\end{proof}

The proof applies to every bounded-degree curve defined over $\Q$
that occurs after a later cut.  Its exceptional alternative uses only the bounded
integral input and output sizes; its Hilbert alternative uses only
degrees and weights.  Thus moving equations and moving integral bases
introduce no new coefficient-height requirement in dimension one.
